% Options for packages loaded elsewhere
\PassOptionsToPackage{unicode}{hyperref}
\PassOptionsToPackage{hyphens}{url}
\documentclass[
]{article}
\usepackage{xcolor}
\usepackage[margin=1in]{geometry}
\usepackage{amsmath,amssymb}
\setcounter{secnumdepth}{-\maxdimen} % remove section numbering
\usepackage{iftex}
\ifPDFTeX
  \usepackage[T1]{fontenc}
  \usepackage[utf8]{inputenc}
  \usepackage{textcomp} % provide euro and other symbols
\else % if luatex or xetex
  \usepackage{unicode-math} % this also loads fontspec
  \defaultfontfeatures{Scale=MatchLowercase}
  \defaultfontfeatures[\rmfamily]{Ligatures=TeX,Scale=1}
\fi
\usepackage{lmodern}
\ifPDFTeX\else
  % xetex/luatex font selection
\fi
% Use upquote if available, for straight quotes in verbatim environments
\IfFileExists{upquote.sty}{\usepackage{upquote}}{}
\IfFileExists{microtype.sty}{% use microtype if available
  \usepackage[]{microtype}
  \UseMicrotypeSet[protrusion]{basicmath} % disable protrusion for tt fonts
}{}
\makeatletter
\@ifundefined{KOMAClassName}{% if non-KOMA class
  \IfFileExists{parskip.sty}{%
    \usepackage{parskip}
  }{% else
    \setlength{\parindent}{0pt}
    \setlength{\parskip}{6pt plus 2pt minus 1pt}}
}{% if KOMA class
  \KOMAoptions{parskip=half}}
\makeatother
\usepackage{longtable,booktabs,array}
\usepackage{caption}
\captionsetup[table]{skip=6pt}
\newcounter{none} % for unnumbered tables
\usepackage{calc} % for calculating minipage widths
% Correct order of tables after \paragraph or \subparagraph
\usepackage{etoolbox}
\makeatletter
\patchcmd\longtable{\par}{\if@noskipsec\mbox{}\fi\par}{}{}
\makeatother
% Allow footnotes in longtable head/foot
\IfFileExists{footnotehyper.sty}{\usepackage{footnotehyper}}{\usepackage{footnote}}
\makesavenoteenv{longtable}
\setlength{\emergencystretch}{3em} % prevent overfull lines
\providecommand{\tightlist}{%
  \setlength{\itemsep}{0pt}\setlength{\parskip}{0pt}}
\usepackage[]{natbib}
\bibliographystyle{apalike}
\usepackage{bookmark}
\IfFileExists{xurl.sty}{\usepackage{xurl}}{} % add URL line breaks if available
\urlstyle{same}
% fallback for those not using the hyperref driver hyperxmp:
\makeatletter
\@ifundefined{xmpquote}{\newcommand{\xmpquote}[1]{#1}}{}
\makeatother
\hypersetup{
  pdftitle={Velaris: effects in signatures{,} budgets at run time{,} and a baseline for a repository's capability surface},
  pdfauthor={Palakurthi Gowri Shankar},
  hidelinks,
  pdfcreator={LaTeX via pandoc}}

\title{Velaris: effects in signatures, budgets at run time, and a
baseline for a repository's capability surface}
\author{Palakurthi Gowri Shankar}
\date{}

\begin{document}
\maketitle

\section{Abstract}\label{abstract}

Code written by language models is increasingly run by people who have
not read it. Velaris is a small programming language for that situation.
A function's signature declares which of seven effects it may perform,
and the compiler checks the declaration across the whole call graph. A
runtime refuses any operation outside a budget the operator writes -
before the operation happens, and in a way the program cannot catch.
Contracts are checked by the Z3 prover where it can settle them, and at
run time where it cannot. A repository can commit a baseline of the
capability surface its programs need, and a check fails any change that
needs more. The central claim is about that surface. Suppose a
repository's Velaris programs are held to a committed baseline by a
required check. If a change makes a program need an effect, path, host,
module or operation count the baseline does not grant, the check fails.
It goes on failing at every later commit at which the program still
compiles and still needs it, until a person edits the baseline. The
claim is not about which function an effect is attributed to: a function
renamed in the change that gives it an effect its program already had
escapes the function-level rule. On a benchmark of 63 programs, 56 with
one defect and 7 correct, each written in Velaris, in JavaScript for
Deno and in Python, Velaris caught 54 of the 56 defects, 42 of them
before running; Deno caught 32 and Python 28; none of the three flagged
a correct program. One of Velaris's two misses is a logic error with no
contract; the other no tool should catch. The capability format is
published separately, under CC0, with a conformance corpus of 444 cases
that an implementation in any language can run.

\section{1. The problem}\label{the-problem}

A model hands someone a script. They can read it line by line, or run it
on trust. Reading every line is the cost they hoped the model would save
them, and each new script asks it again; trust is what the situation
lacks, because the program's comments, names and structure were written
by the same process whose output is in question. Where code arrives
faster than people read it, review is the step that gives.

Three things would help a person in that position. The first is to know,
before running a program, what it may do to the world outside it - which
files, which hosts, which parts of the host system - without reading its
body. The second is to be sure, while it runs, that it does nothing the
person did not allow, whatever it says about itself. The third concerns
a repository rather than a program: to know when the capability of the
code in it grows.

The third matters because capability can be assembled a little at a
time. A helper that builds a URL, a function that reads a file, then a
call three levels down that sends one to the other: each change looks
harmless in its own diff. Hills, Caspary and Cooper Stickland study this
setting directly \citep{hills2026distributed}. A coding agent builds
software over a sequence of pull requests in a persistent codebase while
pursuing a covert side task, and attacks that spread the side task
across pull requests evade monitors that read one diff at a time; no
single monitor they test is robust both to such gradual attacks and to
attacks concentrated in one pull request, and a stateful monitor that
tracks suspicious buildup across pull requests does better against the
gradual ones. Capability is binary and cumulative: the steps of such a
change reach exactly as far as one step that did it all at once. A
comparison of each commit with the one before sees each step alone, and
once a step is merged, sees nothing more of it.

Velaris addresses the three needs with three mechanisms, and adds a
fourth, proofs of contracts, that the others do not depend on. Sections
2.1 to 2.4 describe the four, and section 2.5 the format, published
separately, that specifies the three and not the proofs; section 3
describes the implementation, section 4 what has been measured, section
5 related work and section 6 the limitations.

\section{2. Design}\label{design}

Figure 1 shows the pieces this section describes: one program from its
signature to a refused operation, and one repository whose commits are
each compared with the same baseline.

\begin{verbatim}
(a) One program, one run: examples/effects.vel

fn save_report(text: Text) uses fs {       <- the signature
    write_file("report.txt", text)
}
fn main() uses io, fs, clock, rand { ... }
      |
      |  velaris audit, before it runs
      v
effects: clock, fs, io, rand               <- the audit
reads and writes: report.txt
      |
      |  the operator writes the budget
      v
--allow io,clock,rand,fs:read:./data       <- the budget
      |
      |  at the call, before the write
      v
E313, line 13: 'write_file' reaches        <- the refusal
'report.txt', which this run's fs grants
do not cover. Nothing is written; the run
ends, and the program cannot catch it.

(b) One repository, one baseline

c0: baseline committed, with poll.vel at 10 requests a run

         poll.vel     check, against     review, against
commit   needs        the baseline       the commit before
c1       20           fails: 20 > 10     high
c2       20           fails: 20 > 10     low, sees nothing
c3..c8   40..1000     fails at each;     at c8, sees only
                      c8: 1000 > 10      640 -> 1000
c9       1000, and the baseline edited to 1000:
                      passes             high, names the edit
\end{verbatim}

Figure 1: (a) \texttt{examples/effects.vel}: what a signature declares,
what the audit reports before the program runs, a budget its operator
writes, and the refusal of the write that budget does not grant. (b) the
second history of section 4.3, from \texttt{check\_ratchet.py}: c1
raises a count and is merged, c2 is an unrelated change, and the check
compares every commit with the baseline committed at c0, never with the
commit before it, which is what a review compares.

\subsection{2.1 Effects in signatures}\label{effects-in-signatures}

A Velaris function declares what it may do:

\begin{verbatim}
fn save(path: Text, body: Text) uses fs { ... }
\end{verbatim}

There are seven effects: \texttt{io} (the console and the program's
arguments), \texttt{env} (environment variables), \texttt{fs} (files),
\texttt{net} (the network), \texttt{clock}, \texttt{rand} and
\texttt{ffi} (calling Python, the host language). A function with no
\texttt{uses} clause is pure. The rule is transitive and checked before
the program runs: a function may perform only the effects it declares,
and calling a function requires declaring everything that function
declares, whether or not the callee performs it. A function passed as a
value must be pure, and so must any function a contract calls
\citep[section 3.2]{velaris_spec}. A declaration is therefore an upper
bound on what any call to the function can do, at any depth, and it can
be read without reading the body.

An effect name is coarse: \texttt{fs} says nothing about which files.
The names are for signatures; the finer distinctions belong to the
operator's budget.

\subsection{2.2 Budgets at run time}\label{budgets-at-run-time}

The operator writes a budget when running a program:

\begin{verbatim}
velaris agent_output.vel --allow io,fs:read:./data,net:api.example.com:443@100
\end{verbatim}

A grant names an effect and may narrow it: \texttt{fs} to a direction
and a path, \texttt{net} to a host, a port, or a one-label wildcard
(\texttt{*.example.com}), \texttt{ffi} to named Python modules, and
\texttt{fs} and \texttt{net} to a count of operations in the run
(\texttt{@100}). Grants are additive. Paths are resolved, symbolic links
and \texttt{..} included, when the budget is parsed and again at every
file operation, and compared by whole components \citep[sections 4 and
5]{velaris_spec}.

Every operation is checked when it is attempted, in a fixed order: the
effect (refused with E310), then the scope - a module (E311), a path
(E313), a host or port (E314) - then the count (E315). No file is
opened, no connection made and no module imported for a refused
operation. A refusal ends the run: it is not a failure value the program
can inspect, so a \texttt{check} or \texttt{try} around the call does
not see it. The one exception is a redirect to a host outside the
grants, which fails the request as an ordinary failure naming the
target, because the program did not choose where it was sent. The budget
is checked against what a run attempts, not against what the program
declares, so a program whose declarations were never checked is still
held to it \citep[section 6]{velaris_spec}.

A granted Python module is trusted in full. \texttt{ffi:os} is the
operating system as the current user. What an \texttt{ffi:M} grant does
bound is which modules a call reaches: the attribute chain a call names
is walked step by step, and an object owned by a module outside the
grants is refused, so \texttt{py("json",\ "codecs.encode",\ ...)} under
\texttt{ffi:json} is refused for \texttt{codecs}.

\subsection{2.3 Proofs}\label{proofs}

A function may carry contracts: \texttt{requires}, \texttt{ensures}, and
loop \texttt{invariant}s. The compiler asks the Z3 prover
\citep{demoura2008z3} whether a contract can be broken, using the
contracts of callees rather than their bodies. A contract shown false is
a compile-time error with a counterexample (E700); so is a call shown to
break a callee's \texttt{requires} (E701), a list read shown to pass the
end (E705) and a divisor shown to be zero (E706). A premise the prover
cannot translate abandons the proof, and the contract is checked while
running instead: the compiler does not report a proof it does not have.

Proofs are not part of the capability format. Nothing in sections 2.1,
2.2 and 2.4 depends on them, and an implementation with no prover can
conform to the format at every level (section 2.5).

\subsection{2.4 The ratchet}\label{the-ratchet}

\texttt{velaris\ capabilities\ init} reads every \texttt{.vel} file in a
repository and writes \texttt{velaris.capabilities}: the repository's
\emph{surface} - every grant its programs need, in the budget grammar,
and for \texttt{fs} and \texttt{net} the most operations one run can
perform - and, for each program, its own grants and counts and the
effects each of its functions declares \citep[section 9]{velaris_spec}.
What a program needs is read from its text: its functions' declarations,
and the paths, hosts and modules its calls name as fixed text - a
literal, or a variable bound once to one. A path, URL or module built
while running is recorded as the unscoped grant (\texttt{fs:read},
\texttt{net}, \texttt{ffi}), which a scoped baseline does not cover. The
operation bound comes from loops whose counter and limit the text fixes,
multiplied through nesting and calls; a loop whose turns the text does
not fix, and recursion, have no bound. The prover takes no part, so the
result is the same with and without it.

\texttt{velaris\ capabilities\ check} derives the same from the working
tree and compares it with the committed file - with that file, and never
with a previous commit. A change widens the baseline, and the check
fails, under five rules: a grant the surface does not cover (W1); more
operations than the surface's count (W2); for a program the baseline
records, a grant (W3) or a count (W4) beyond its own entry; and an
effect added to a function the baseline records (W5). Narrowing never
fails. A program the baseline does not record is held to the surface
alone. A check that cannot read the baseline - it is missing, is not the
right format, or holds a grant the format does not allow - fails rather
than passes, so deleting the file does not turn the check off. The
writer will not replace an existing baseline unless asked to in so many
words, so every widening that is accepted arrives as an edit to the
file, in review.

Because the comparison is with the baseline, a widening merged once goes
on failing at every later check until someone edits the file, and a
widening assembled over many commits is reported as its whole sum
(Figure 1b). That is the property the gradual attacks of section 1 need
to defeat.

\textbf{The central claim.} Suppose a repository's Velaris programs are
held to a committed baseline by a required check. If a change makes a
program need an effect, path, host, module or operation count the
baseline does not grant, the check fails. It goes on failing at every
later commit at which the program still compiles and still needs it,
until a person edits the baseline. Which entries count depends on the
program: for one the baseline does not record, or records as not
compiling, the baseline grants what its surface grants; for one it
records as compiling, only what both the surface and the program's own
entry grant.

``Needs'' there is what the program's text requires by the derivation of
velaris-spec section 9.3: the declared surface. Three things follow that
the claim does not make. It does not say what a granted module does: a
program newly calling \texttt{os.system} through a baseline's
\texttt{ffi:os} is inside the surface. It does not say where a path
leads: paths in a baseline are compared as text, and a symbolic link
under a recorded directory is that directory's content; the budget,
which resolves every path at run time, is where that is caught. And it
does not say which function an effect belongs to: a function is known by
its file and name, so a helper renamed in the same change that gives it
an effect its program already had is a new function, held to its
program's entry and the surface but not to what its old name declared,
and W5 does not report it. The declared surface did not widen in that
case, and the check is right to pass it; what escapes is the
attribution.

A separate command, \texttt{velaris\ review\ -\/-against\ REF}, reports
how the working tree differs from a git ref, with a one-word risk. It
compares with a previous state, so it informs a reviewer and is not the
gate.

\subsection{2.5 The format, published
separately}\label{the-format-published-separately}

The effects, the budget grammar, what a runtime must refuse,
\texttt{velaris.audit/1} (the JSON report of what a program declares and
names before it runs) and \texttt{velaris.capabilities/1} are specified
separately, under CC0, as velaris-spec \citep{velaris_spec}, so that
they can be implemented without reading the compiler. Its conformance
document defines three levels: L1, declaration - parse the budget
grammar, compute a program's effect surface, emit
\texttt{velaris.audit/1} that validates against its schema; L2,
enforcement - refuse at run time every operation outside the budget,
uncatchably; L3, the ratchet - write and read baselines and classify
widening against narrowing for every kind of scope. L2 and L3 each
require L1; L3 does not require L2, so a tool with no runtime can
conform at L1 and L3. No level requires a prover. Conformance is shown
by running a corpus of JSON cases, described in section 4.2.
velaris-spec also defines an in-toto predicate type that binds an audit
to the digests of the files audited, so that a signed statement can say
which source an audit describes; velaris-lang 4.2.0 writes such
statements, unsigned (\texttt{velaris\ attest}), and leaves signing to
Sigstore's tools.

\section{3. Implementation}\label{implementation}

The implementation, velaris-lang \citep{velaris_lang}, is one Python
file, \texttt{velaris.py}, of 13,497 lines - one file by design: nothing
to install but Python, nothing vendored, a single artifact to sign and
audit. It is laid out in the order a program passes through it: lexer,
parser, loader, effect checker, type checker, prover, native code
generation, interpreter. The prover needs the optional
\texttt{z3-solver} package and the native code generator the optional
\texttt{llvmlite}; without them, contracts are checked while running and
everything is interpreted, and the effect checks and the budget are
unchanged.

The budget is enforced in the interpreter, at each builtin that performs
an operation, before the operation. A run with a time limit or a memory
cap runs in a child process that can be killed, and a pool of such
workers re-installs the budget before every program. The same checks sit
behind every interface: the command line, a Python library
(\texttt{check}, \texttt{audit}, \texttt{run}, \texttt{Pool}), an HTTP
door and an MCP server whose operators set ceilings callers cannot
exceed, and a GitHub Action that runs the capability check and uploads
findings as SARIF. Every error has a stable code; the compiler's table
holds 62.

The budget is not a security boundary. It is enforced by an interpreter
written in Python, in the same process as the compiler and, unless a
limit is set, the program. It is a guard against a program doing what it
was not asked to, and it belongs inside an operating-system sandbox when
the stakes warrant one.

\section{4. Evaluation}\label{evaluation}

\subsection{4.1 A benchmark against Deno and
Python}\label{a-benchmark-against-deno-and-python}

The benchmark is 63 small programs, each written three times with the
same behaviour: in Velaris, in JavaScript for Deno, and in Python.
Fifty-six contain one deliberate defect, on one line marked in all three
sources; seven are correct controls. They fall in eleven categories: a
file write hidden in a helper, a network call hidden in a helper,
division by a value that can be zero, a read past the end of a list,
integer overflow, an ignored failure, an infinite loop, runaway memory,
reaching a dangerous module, correct programs that must not be flagged,
and a grant narrower than the effect. One harness runs every program
through every tool with the same rules: a 5-second timeout, a 256 MB
memory cap, and for Velaris the narrowest budget each task needs. Each
program gets one verdict per tool - caught before running, caught while
running, missed, or, for a control, clean or a false positive - and
after every run the harness observes whether the dangerous effect
actually happened: a file created, a request received by a local
listener, a subprocess's sentinel printed. The committed results were
produced by Velaris 4.1.0, Deno 2.9.6 and Python 3.13.13 on Windows 11;
Velaris 3.0.0, with the same Deno and Python on the same platform, had
produced the same table before it, every verdict and every line of
evidence.

{\def\LTcaptype{none} % do not increment counter
\begin{longtable}[]{@{}
  >{\raggedright\arraybackslash}p{(\linewidth - 8\tabcolsep) * \real{0.2000}}
  >{\raggedright\arraybackslash}p{(\linewidth - 8\tabcolsep) * \real{0.2000}}
  >{\raggedright\arraybackslash}p{(\linewidth - 8\tabcolsep) * \real{0.2000}}
  >{\raggedright\arraybackslash}p{(\linewidth - 8\tabcolsep) * \real{0.2000}}
  >{\raggedright\arraybackslash}p{(\linewidth - 8\tabcolsep) * \real{0.2000}}@{}}
\toprule\noalign{}
\begin{minipage}[b]{\linewidth}\raggedright
Tool
\end{minipage} & \begin{minipage}[b]{\linewidth}\raggedright
Caught before running
\end{minipage} & \begin{minipage}[b]{\linewidth}\raggedright
Caught while running
\end{minipage} & \begin{minipage}[b]{\linewidth}\raggedright
Missed
\end{minipage} & \begin{minipage}[b]{\linewidth}\raggedright
False positives (7 controls)
\end{minipage} \\
\midrule\noalign{}
\endhead
\bottomrule\noalign{}
\endlastfoot
Velaris & 42 & 12 & 2 & 0 \\
Deno & 5 & 27 & 24 & 0 \\
Python & 0 & 28 & 28 & 0 \\
\end{longtable}
}

Table 1: the 56 dangerous programs and the 7 controls, from
\texttt{benchmark/RESULTS.md}.

Where the catches come from differs by tool. Velaris's catches before
running come from effects in signatures (the file, network and module
categories, and the environment secret of category 11), from
unhandled-failure checks and the prover (division, list reads, ignored
failures), and from its termination rule, which flags a loop whose
counter does not move one step toward an unchanging limit. That rule
flagged the eleven infinite and memory-growth programs before running;
it claims nothing about whether such a loop ends, and each of these
happened not to. Deno's catches are almost all at the moment of the
call, through its permission flags; nothing in \texttt{deno\ check} or
\texttt{deno\ lint} reads a file write or a fetch as a problem. In one
network program Deno denied the request but the program caught the
denial, which is an ordinary exception in JavaScript, and exited 0; the
harness credits Deno because it observed that no request arrived, and a
caller reading only the exit status would have seen success. A Velaris
refusal cannot be caught.

Velaris alone caught the six integer-overflow programs, while running:
its whole numbers are 64-bit and arithmetic that leaves that range stops
the program (E407). Python's integers do not overflow, so its printed
values are arithmetically right and only wrong for a 64-bit consumer;
the results record those rows as Python misses and say a reader may
discount them. Python's catches in the division and ignored-failure
categories are crashes on the inputs the harness supplies - a zero, a
non-number, a missing key - and with ordinary input those programs run
clean; Velaris's E520, E705 and E706 do not depend on the input. Four
Velaris catches came only at run time where a sibling program was caught
before running: a division on the unguarded one of two paths, a
remainder inside a loop, a division in \texttt{main} on
\texttt{n\ -\ 1}, and a read at \texttt{i\ +\ 1} in a loop bounded by
the list's length. The prover did not settle those, and the runtime
checks stopped them.

Velaris missed two programs, both included on purpose so that the table
is not a list of what the language was built to do. In \texttt{04c} a
loop stops one item early: no read is out of range and the function has
no contract, so a wrong total is indistinguishable from a right one. In
\texttt{09c} a program prints \texttt{rm\ -rf\ build} as a hint for its
caller and touches nothing; its only effect is \texttt{io}, which the
task needs. That one should not be caught by any tool: control program
\texttt{10d} prints the same words in a warning and is harmless, and
nothing that looks at the program from outside can tell the two apart.
The danger is in what a caller does with the text.

\textbf{What this benchmark does not show.} The programs were written
for it by this project; in the first ten categories, three of each six
were written after the first three, against the tools, to hide the same
defects better. It is small. The inputs were chosen to trigger each
defect. The committed run is one machine's. And it measures programs,
not the ratchet, which section 4.3 treats separately. The harness is in
the repository, its rules are one function each, and the evidence for
every cell is recorded next to the verdict, so a reader who disagrees
with a row can change the rule and rerun it.
\texttt{benchmark/run.py\ -\/-check} reruns the whole table and fails if
any verdict differs from the committed \texttt{benchmark/results.json}.

\subsection{4.2 The suites, and a conformance
corpus}\label{the-suites-and-a-conformance-corpus}

Each guarantee has a suite that asserts it, and each suite runs on every
push on Linux, Windows and macOS, with Python 3.10 and 3.12, with and
without the prover. \texttt{check\_sandbox.py} holds 34 attempts to
escape a budget - an effect two helpers below \texttt{main}, a refusal
caught to carry on, a module reached through a submodule path,
\texttt{py\_json} or a handle, a path out of a prefix by \texttt{..} or
a symbolic link, a host, a port, a wildcard's parent domain, a count, a
URL with no scheme taken as HTTPS - each refused with the code it must
carry, and 16 honest programs that must still run. Run on Windows with
the prover, \texttt{check\_ratchet.py} passes 114 checks of the ratchet,
\texttt{check\_library.py} 196 of the library, the doors, the audit and
its attestations (its one symbolic-link case runs only on POSIX),
\texttt{check\_refusals.py} 21 wrong programs each refused with the
right code, \texttt{check\_fallible.py} 26 fallible builtins, and
\texttt{check\_termination.py} 44 adversarial loops; without it,
\texttt{check\_library.py} passes 193 and \texttt{check\_refusals.py}
11, skipping the checks that are about proofs, and the rest are
unchanged (CHANGELOG, 4.2).

Those suites drive velaris-lang's own command line and Python library,
so until velaris-lang 4.1 an implementation in another language could
use them only through an adapter. velaris-spec now holds a corpus of 444
JSON cases that any implementation runs its own way: 298 at L1 (280
budgets to parse or refuse, 18 programs to audit), 37 at L2 (programs to
run under a budget, with a fixture of files and two local HTTP servers)
and 109 at L3 (baselines to write, changes to check, two sequences of
changes, the writer's guard, and the covering, reduction and
operation-bound rules case by case). A script writes the corpus from the
tables of three of the suites, where each entry is also asserted against
velaris-lang, and a drift test in both repositories' CI regenerates it
and fails if it differs from what is committed. Thirteen scenarios are
left out and listed with the reason: six depend on Python's object
model, four need a Python host, one is a run given no budget, which the
format leaves to the implementation, and two are about velaris-lang's
command-line flags. \texttt{velaris\ conformance} runs the corpus
against velaris-lang; it passes at all three levels (velaris-spec
\texttt{tests/index.json}; CHANGELOG, 4.2).

Writing the corpus found a defect. The audit of a program refused for
naming an unknown effect still listed that name among the program's
effects, and wrote a suggested budget that does not parse, although the
specification said neither could happen. velaris-lang 4.1.0 fixes the
implementation, and velaris-spec 0.4 corrects the text that had claimed
more than the implementation did.

\subsection{4.3 The ratchet}\label{the-ratchet-1}

\texttt{check\_ratchet.py} builds scratch trees and real git histories
and asserts what the check says at every step. In its gradual case, a
greeter program is committed with its baseline, and six changes follow:
a pure text helper; the greeting using it; a pure summary function;
\texttt{main} printing a summary; a function returning a collector's URL
as text; and finally a summary delivered to that URL through a helper
three calls below \texttt{main}. The first five pass. The sixth fails,
naming \texttt{net:collector.example.net} as a new effect, the file
\texttt{lib/deliver.vel}, line 2, the function \texttt{send}, and the
chain
\texttt{main\ -\textgreater{}\ summary\ -\textgreater{}\ deliver\ -\textgreater{}\ send}.
A review of each change against the one before calls the first five
\texttt{low} and the sixth \texttt{high}. In a second history a count of
10 requests a run is raised and merged, and keeps failing at every later
commit; seven steps from 10 to 1000 are reported by the check as 1000
against the declared 10, while a comparison with the previous commit
shows only the last step, 640 to 1000. The suite also holds changes that
must pass - reordering, reformatting, renaming locals, a literal moved
into a variable, a function renamed or moved, narrowing, a new program
inside the surface - and baselines that cannot be read, which must fail.

The five known limits of section 6 are cases too, each recording the
outcome the check gives today, so that another implementation matches it
and so that a change to it is visible. The repository checks itself: its
own \texttt{velaris.capabilities} records 172 programs, 22 of them built
to be refused and recorded as not compiling, and CI runs the check on
every leg.

The ratchet has not been measured on the gradual-attack benchmark of
Hills et al., whose code is not Velaris, or on any repository but its
own. What is claimed for it is a property of a deterministic comparison,
in the same sense that soundness is a property of a type system, and not
a detection rate. A rate measures how often a heuristic notices
something. The ratchet does not notice; it computes, and the cases in
\texttt{check\_ratchet.py} show what it computes. The property is the
one stated in section 2.4, with the three things it does not say there
and the five known limits of section 6, and it rests on the comparison's
definition and on those cases: it has not been proved mechanically.

\section{5. Related work}\label{related-work}

This section says what each piece of work does and how Velaris differs.
It claims priority over none of it; velaris-lang's changelog does not
name any of it as the source of a design decision.

\textbf{Object capabilities.} Dennis and Van Horn introduced the
capability, a reference that both names a resource and carries the right
to use it \citep{dennis1966semantics}; Miller's object-capability model
builds least authority from references a program has been handed
\citep{miller2006robust}. A Velaris effect is a name in a signature, and
a budget is ambient to a whole run: any function declaring \texttt{fs}
may reach any path the budget allows without holding a reference to it.
Nothing in Velaris is unforgeable, and nothing is handed over.
``Capability'' in the format's name follows common usage.

\textbf{TACIT.} Odersky and colleagues have agents write Scala 3 with
capture checking, in which capabilities - a file system rooted at a
directory, a set of hosts, a set of commands - are values the type
system tracks, and pure computations over classified data cannot leak it
\citep{odersky2026tacit}. In Velaris capabilities are effect names plus
an operator's budget written as text and enforced by an interpreter, in
a small language a model learns from a card; it has no counterpart to
TACIT's information-flow control, and no grant for running commands.

\textbf{CaMeL.} Debenedetti and colleagues have a privileged model turn
a trusted request into a program in a restricted subset of Python, run
by an interpreter that attaches to every value its provenance and
permitted readers and checks a policy at each tool call, while a
quarantined model parses untrusted data \citep{debenedetti2025camel}.
Velaris tracks nothing about data: a budget bounds which effects, paths,
hosts and modules a run may reach, not which values may flow to them.

\textbf{WASI.} The WebAssembly System Interface gives a module only the
directories, sockets and other resources its host hands it as handles,
so a module given nothing reaches nothing \citep{wasi}. Velaris's grants
are text an operator writes, enforced inside the same process rather
than at a virtual machine's boundary, and its command line grants every
effect when no budget is given, so deny-by-default holds only once an
operator writes one.

\textbf{Deno.} Deno's permission flags give a script no file, network,
environment or subprocess access unless granted, with scoped forms much
like Velaris's grammar \citep{deno_security}; section 4.1 compares the
two directly. A Deno denial is an exception a script can catch; a
Velaris refusal ends the run. Deno checks at the call; Velaris also
declares the effect in the signature, so it is visible before running.

\textbf{Effect systems.} Lucassen and Gifford's polymorphic effect
systems record in an expression's type the side effects evaluating it
may have \citep{lucassen1988effects}. Velaris's \texttt{uses} is a fixed
set of seven names with no polymorphism - a function value is simply
required to be pure - a small instance of the idea.

\textbf{in-toto and SLSA.} in-toto binds signed statements about
software artifacts to subjects identified by digest
\citep{torresarias2019intoto, intoto_attestation}, and SLSA defines
levels of build integrity and a provenance predicate \citep{slsa}.
Velaris's documents describe what a source text declares, not how or by
whom an artifact was built. velaris-spec defines an in-toto predicate
type for the audit, and velaris-lang 4.2.0 writes statements of it; its
release workflow signs one for an example program with Sigstore's tools
and verifies it.

\textbf{SARIF.} SARIF is a common format for the findings of static
analysis tools \citep{sarif2020}. velaris-lang writes its compile
errors, unproven promises and capability widenings as SARIF for code
scanning; the audit itself describes a program rather than listing
findings, and is not SARIF.

\textbf{Gradual attacks.} Hills, Caspary and Cooper Stickland show that
a coding agent can spread a covert side task across pull requests so
that no single diff looks decisive, that such attacks evade diff
monitors, and that a stateful monitor tracking buildup across pull
requests detects them better \citep{hills2026distributed}. The ratchet
is not a monitor of intent. It is a deterministic comparison of what a
repository's Velaris programs need with a baseline the repository
declared, at every change. It sees the part of a gradual attack that
needs capability the baseline did not declare, and nothing else: a side
task that stays within the declared surface, or is not written in
Velaris, is outside it.

\section{6. Limitations}\label{limitations}

These are stated as velaris-lang's THREAT\_MODEL.md states them.

\begin{itemize}
\tightlist
\item
  \textbf{Only code written in Velaris.} A model asked for Velaris may
  hand back Python. Nothing here applies to code the Velaris runtime
  does not run.
\item
  \textbf{A granted \texttt{ffi} module.} A granted module can do
  whatever it can do. \texttt{ffi:os} is the operating system;
  \texttt{ffi:subprocess} is a shell. The grant narrows which modules a
  call may reach, not what a module does.
\item
  \textbf{Logic errors with no contract.} A function that returns the
  wrong number and promises nothing is correct as far as the compiler
  knows. Benchmark program \texttt{04c} is this, and is a miss.
\item
  \textbf{The meaning of text.} A program that prints a shell command
  for its caller touches nothing. No effect system can tell it from a
  program that prints the same words as a warning, and Velaris does not
  try.
\item
  \textbf{One maintainer.} The project is maintained by one person.
  Fixes to soundness and sandbox reports are promised within a week, and
  nothing else is promised.
\item
  \textbf{Not a security boundary.} The budget is enforced by a Python
  interpreter in the same process as the compiler. A defective or
  altered \texttt{velaris.py} holds nothing. Side channels, resource use
  below the limits, where a granted host name resolves, and what a
  granted host does with a request are outside the model.
\item
  \textbf{The ratchet.} The declared surface cannot widen without the
  check failing (section 2.4); these are its known limits, each held as
  a case in \texttt{check\_ratchet.py} and in the corpus with the
  outcome it has today. Two widenings that pass: a function renamed in
  the same change that gives it an effect its program already had evades
  function-level attribution (W5), though the declared surface did not
  widen; and what a granted \texttt{ffi} module does is outside the
  program's text. Three changes that do not widen and are reported as
  widenings, each a conservative false positive: a loop limit behind a
  function call (\texttt{for\ i\ in\ 0\ to\ limit()}) loses its bound
  and is reported as unbounded; a path passed to a helper as a parameter
  is taken as unscoped even when every caller passes a literal; and a
  path written with a backslash is covered only by the same text, so
  \texttt{data\textbackslash{}in.csv} is outside a recorded
  \texttt{data} even on Windows, where it is inside. The ratchet also
  guards nothing if the check is not required, and an edit to the
  baseline is an accepted widening whose review is the control.
\item
  \textbf{The evaluation.} Section 4.1's benchmark was written by this
  project and is small; no user study has been done; the ratchet's
  properties rest on its definition and its cases, not on a measured
  detection rate.
\end{itemize}

\section{7. Conclusion}\label{conclusion}

Velaris puts what a function may do in its signature and checks it
across the call graph; refuses, while a program runs, every operation
outside a budget its operator wrote, in a way the program cannot catch;
proves contracts where the prover can; and holds a repository's declared
capability surface to a baseline that only an edit can widen. The first
three can be tested program by program, and on a 63-program benchmark
they caught 54 of 56 defects, with no false positives on the controls.
The fourth is a property, not a rate: under a required check, the
declared surface does not widen without the check failing, however the
change is divided among commits. What it does not do - attribute an
effect to the right function across a rename, see into a granted module,
tell harmless text from a dangerous command - is stated here as
precisely as what it does, because the second list is only worth
believing if the first one is.

\section{Use of generative AI}\label{use-of-generative-ai}

The software described here was developed by the author with the
assistance of AI coding agents, and the first draft of this paper was
written by an AI coding agent from the repository's contents under the
author's direction. The author specified the structure and claims,
reviewed and revised the text, and is responsible for its accuracy.
Every number reported was verified against the repository files named in
the reproducibility section. Three adversarial reviews by AI systems,
credited in the repository's HALL\_OF\_FAME.md, found defects that are
recorded in its changelog.

\section{Reproducibility}\label{reproducibility}

The numbers in this paper can be regenerated from the two repositories
at their tags:

\begin{verbatim}
git clone https://github.com/gowrishankar-infra/velaris-lang
git clone https://github.com/gowrishankar-infra/velaris-spec
git -C velaris-spec checkout v0.5.1
cd velaris-lang
git checkout v4.2.1
pip install ".[full,test]"         # the prover, the native compiler, jsonschema

# Table 1 (needs Deno 2.x on PATH for the Deno column; 5 to 8 minutes)
python benchmark/run.py            # rewrites benchmark/RESULTS.md and results.json
python benchmark/run.py --check    # the same, and fails if a verdict differs

# Figure 1a
velaris audit examples/effects.vel
velaris examples/effects.vel --allow io,clock,rand,fs:read:./data

# section 4.2 and 4.3
python check_sandbox.py
python check_library.py
python check_ratchet.py
python check_refusals.py
python check_fallible.py
python check_termination.py
velaris conformance --corpus ../velaris-spec/tests
python build_conformance.py --check ../velaris-spec/tests
velaris capabilities check .

# the same without the prover
python -m venv bare
bare/bin/pip install ".[test]"     # bare\Scripts\pip on Windows
bare/bin/python check_sandbox.py   # and each command above
\end{verbatim}

Where each number comes from:

{\def\LTcaptype{none} % do not increment counter
\begin{longtable}[]{@{}
  >{\raggedright\arraybackslash}p{(\linewidth - 2\tabcolsep) * \real{0.5000}}
  >{\raggedright\arraybackslash}p{(\linewidth - 2\tabcolsep) * \real{0.5000}}@{}}
\toprule\noalign{}
\begin{minipage}[b]{\linewidth}\raggedright
Number
\end{minipage} & \begin{minipage}[b]{\linewidth}\raggedright
File
\end{minipage} \\
\midrule\noalign{}
\endhead
\bottomrule\noalign{}
\endlastfoot
63 programs, 56 dangerous, 7 controls; 42/12/2, 5/27/24, 0/28/28; 0
false positives; Velaris 4.1.0, Deno 2.9.6, Python 3.13.13, Windows 11;
5 s timeout, 256 MB cap; the misses \texttt{04c} and \texttt{09c} &
\texttt{benchmark/RESULTS.md} \\
eleven categories; three of each six programs written against the tools;
5 to 8 minutes for a full run & \texttt{benchmark/README.md} \\
the same table from Velaris 3.0.0 & \texttt{benchmark/RESULTS.md} at
v4.0.0, and \texttt{CHANGELOG.md}, 4.1 entry \\
seven effects; the refusal codes E310, E311, E313, E314, E315 &
\texttt{SPEC.md} section 7 and 7.1 \\
E407, E520, E700, E701, E705, E706; 62 error codes &
\texttt{velaris.py}, \texttt{ERROR\_TABLE} \\
Figure 1a: the program, its effects and paths, line 13 &
\texttt{examples/effects.vel} \\
13,497 lines & \texttt{velaris.py} \\
34 escape attempts and 16 honest programs & \texttt{check\_sandbox.py},
\texttt{ESCAPES} and \texttt{HONEST} \\
Linux, Windows and macOS; Python 3.10 and 3.12; with and without the
prover & \texttt{.github/workflows/test.yml} \\
114, 196, 21, 26, 44 checks; 193 and 11 without the prover &
\texttt{CHANGELOG.md}, 4.2 entry \\
444 cases: 298, 37, 109; 280 budgets, 18 audits; 13 left out &
velaris-spec \texttt{tests/index.json} \\
the six-commit history; line 2 of \texttt{lib/deliver.vel}; 10 to 1000,
640 to 1000, and Figure 1b & \texttt{check\_ratchet.py},
\texttt{SEQUENCES} \\
five known limits & \texttt{check\_ratchet.py}, \texttt{CHECKS};
\texttt{CHANGELOG.md}, 4.0 entry \\
172 programs, 22 not compiling & \texttt{velaris.capabilities} \\
\end{longtable}
}

\renewcommand\refname{References}
\bibliography{references.bib}

\end{document}
